
\section{\BookRecordIndex{T0246}模空间：\BookRecordIndex{T0408}正则性与紧致性}\label{sec:1-2}
\zhindex{mofangjian}{模空间}
\zhindex{zhengzexing}{正则性}
\zhindex{jinzhixing}{紧致性}

\BookRecordIndex{T0003}$J$-全纯曲线的关键事实是：当 $J$ 处于一般位置时，这些曲线组成有限维族。这些族构成
有限维流形
\[
  \mathcal M^*(A;J),
\]
称为\BookRecordIndex{T0246}模空间；其配边类与在 $\mathcal J_\tau(M,\omega)$ 中具体选择哪个 $J$ 无关。
这里 $A\in H_2(M,\mathbb Z)$ 是一个同调类，而 $\mathcal M^*(A;J)$ 基本上由所有表示
$A$ 的 \BookRecordIndex{T0003}$J$-全纯曲线 $u:\Sigma\to M$ 组成。虽然流形 $\mathcal M^*(A;J)$ 几乎从不
紧致，但它通常保留了足够的紧致性，使我们能够用它定义不变量。

本书第 2--5 章用来精确表述并证明上述论断。结果自然分为两组：第一组讨论\BookRecordIndex{T0408}正则性和
\BookRecordIndex{T0166}横截性问题，见第 2、3 章；第二组讨论 $\mathcal M^*(A;J)$ 的紧化，见第 4、5 章。
下面几段选出最重要的概念和定理。初学这一理论的读者首先应把注意力放在这些结果及其证明
上。每章的导言还会进一步说明本书的阅读方法。

\paragraph{局部性质。}
第 2 章讨论 \BookRecordIndex{T0003}$J$-全纯曲线的局部性质。对后续发展而言，关键结果是\BookRecordIndex{T0249}能量恒等式和命题
2.5.1：前者是后面紧致性结果的基础，后者刻画了哪些曲线不是多重覆盖的。若曲线
$u:\Sigma\to M$ 可写成一个次数大于 $1$ 的\BookRecordIndex{T0287}全纯分歧覆盖映射
$(\Sigma,j)\to(\Sigma',j')$ 与一个 $J$-全纯映射 $\Sigma'\to M$ 的复合，就称 $u$
是多重覆盖的；不是多重覆盖的曲线称为\BookRecordIndex{T0091}单纯曲线。\BookRecordIndex{T0113}多重覆盖曲线常常是由所有
$J$-全纯 $A$-曲线组成的\BookRecordIndex{T0246}模空间 $\mathcal M(A;J)$ 的奇点，因此需要一种可操作的判据
来保证 $u$ 单纯。如果存在 $z\in\Sigma$ 使得
\[
  du(z)\ne0,\qquad u^{-1}(u(z))=\{z\},
\]
就称曲线 $u$ \BookRecordIndex{T0397}在某处单射；具有这一性质的点 $z$ 称为 $u$ 的\BookRecordIndex{T0094}单射点。命题 2.5.1 的
内容如下。它是定理 3.1.6 证明中的重要一环；后者断言，对一般的被 $\omega$ \BookRecordIndex{T0372}驯服的
$J$，单纯 $J$-全纯 $A$-曲线的\BookRecordIndex{T0246}模空间 $\mathcal M^*(A;J)$ 是光滑流形。

\begin{namedproposition}
每条单纯 \BookRecordIndex{T0003}$J$-全纯曲线 $u$ 都\BookRecordIndex{T0397}在某处单射。此外，\BookRecordIndex{T0094}单射点的集合在 $\Sigma$ 中既开又稠密。
\end{namedproposition}

除了一处要引用一个把\BookRecordIndex{T0333}唯一延拓原理推广到 \BookRecordIndex{T0003}$J$-全纯曲线的深刻分析定理外，证明是初等的。

\paragraph{弗雷德霍姆理论。}
固定一个亏格为 $g$ 的黎曼曲面 $\Sigma$，以 $\mathcal M^*(A;J)$ 表示所有代表类 $A$
的单纯 $J$-全纯映射 $u:\Sigma\to M$ 组成的集合。
\footnote{本书第一版把这个空间简单记作 $\mathcal M(A;J)$。不过，如今惯常用字母
$\mathcal M$ 表示完整的\BookRecordIndex{T0246}模空间。}
下面是第 3 章主要结果的非正式版本。回忆：完备度量空间的一个子集，如果包含可数多个
开稠密集的交，就称为（贝尔意义下的）\BookRecordIndex{T0300}剩余集。贝尔纲定理断言，每个这样的集合都是稠密的。

\begin{namedtheorem}[A]
存在剩余子集
$\mathcal J_{\mathrm{reg}}(A)\subset\mathcal J_\tau(M,\omega)$，使得对每个
$J\in\mathcal J_{\mathrm{reg}}(A)$，空间 $\mathcal M^*(A;J)$ 都是维数为
\[
  \dim\mathcal M^*(A;J)=n(2-2g)+2c_1(A)
\]
的光滑流形。此流形 $\mathcal M^*(A;J)$ 带有一个自然定向。
\end{namedtheorem}

另一个重要定理刻画 $\mathcal M^*(A;J)$ 如何依赖于 $J$ 的选择（定理 3.1.8）。这些
定理成立的根本原因是\BookRecordIndex{T0209}柯西--黎曼方程
\[
  du\circ j=J\circ du
\]
是椭圆型的，因而它在一条 \BookRecordIndex{T0003}$J$-全纯曲线处的线性化是一个\BookRecordIndex{T0128}弗雷德霍姆算子，记为 $D_u$。
定理 A 中的集合 $\mathcal J_{\mathrm{reg}}(A)$，正是所有满足下述条件的\BookRecordIndex{T0088}殆复结构
$J\in\mathcal J_\tau(M,\omega)$ 组成的集合：对每个 $u\in\mathcal M^*(A;J)$，
\BookRecordIndex{T0346}线性化算子 $D_u$ 都是满射。$\mathcal J_{\mathrm{reg}}(A)$ 的元素有时称为正则\BookRecordIndex{T0088}殆复结构。
一个有趣的事实是，$J$ 的\BookRecordIndex{T0372}驯服条件在这里无关紧要；上述结果对任意紧流形 $M$ 上的所有
\BookRecordIndex{T0088}殆复结构 $J$ 都成立。

巴拿赫空间之间的有界线性算子 $D:X\to Y$，若其核有限维，且像闭并在 $Y$ 中具有
有限余维，就称为\BookRecordIndex{T0128}弗雷德霍姆算子。$D$ 的指标定义为核与余核维数之差：
\[
  \operatorname{index}D=\dim\ker D-\dim\operatorname{coker}D.
\]
一个重要事实是，在范数拓扑下，\BookRecordIndex{T0128}弗雷德霍姆算子的集合是开集，而且\BookRecordIndex{T0129}弗雷德霍姆指标在该
集合的每个连通分支上为常数。因此，当 $D$ 连续变化时指标不变，尽管核和余核的维数当然
都可能改变。

我们在附录 A 中说明，\BookRecordIndex{T0128}弗雷德霍姆算子本质上与有限维算子一样具有良好性质，并且在
无限维\BookRecordIndex{T0380}隐函数定理中起重要作用。更准确地，设 $F:X\to Y$ 是\BookRecordIndex{T0213}可分巴拿赫空间之间的
$C^\infty$ 光滑映射，并且在每一点 $x\in X$，其微分 $dF(x):X\to Y$ 都是指标为
$k$ 的\BookRecordIndex{T0128}弗雷德霍姆算子。若 $y\in Y$ 是 $F$ 的\BookRecordIndex{T0409}正则值，即对每个 $x\in F^{-1}(y)$，
$dF(x)$ 都为满射，那么与有限维情形一样，逆像
\[
  \mathcal M:=F^{-1}(y)
\]
是一个 $k$ 维光滑流形。\BookRecordIndex{T0296}萨德定理的无限维版本断言，$Y$ 中几乎所有点都是 $F$ 的\BookRecordIndex{T0409}正则值
（严格说来，\BookRecordIndex{T0409}正则值组成\BookRecordIndex{T0300}剩余集）。当 $X,Y$ 是\BookRecordIndex{T0213}可分\BookRecordIndex{T0053}巴拿赫流形而非巴拿赫空间时，这个定理
仍然成立。然而，按上述形式，它不能推广到其他类型的无限维向量空间，例如弗雷歇空间。
所以，尽管定理 A 中的集合 $\mathcal J_{\mathrm{reg}}$ 确实由一个以
$\mathcal J_\tau(M,\omega)$ 为值域的\BookRecordIndex{T0128}弗雷德霍姆算子的\BookRecordIndex{T0409}正则值组成，证明中仍有一些额外
技术问题，因为 $\mathcal J_\tau(M,\omega)$ 是\BookRecordIndex{T0130}弗雷歇流形而不是\BookRecordIndex{T0053}巴拿赫流形。

\paragraph{$J$-全纯圆盘。}
尽管本书大多数结果讨论 $J$-全纯球面的性质，我们在第 3、4 章以及附录 B、C 中，是对
带有可能为空的边界的一般紧黎曼曲面表述并证明基础结果的，因而也包括 $J$-全纯圆盘。
第 3 章将证明，在拉格朗日边界条件下，$\bar\partial$ 算子仍是弗雷德霍姆的。因此，
边界落在给定\BookRecordIndex{T0219}拉格朗日子流形中的 $J$-全纯圆盘之\BookRecordIndex{T0246}模空间，对一般的 $J$ 仍然构成有限维
流形。这一事实是附录 C 所给\BookRecordIndex{T0221}黎曼--罗赫定理归纳证明中的重要要素。第~\ref{sec:9-2} 节给出这一理论
的若干应用，尤其是证明 $\mathbb R^{2n}$ 中不存在闭的恰当\BookRecordIndex{T0219}拉格朗日子流形。拉格朗日边界
条件是椭圆型的这一事实，在弗洛尔理论中也非常重要；参见第 12 章。

这是把理论推广到具有更一般定义域的映射的第一步。下一步可以考虑亏格较高的闭定义域，
或者亏格为零的非紧定义域。本书的工作足以处理复结构 $j$ 固定的高亏格闭定义域
$(\Sigma,j)$。改变 $j$ 不会显著影响弗雷德霍姆理论，却会在讨论紧致性时引出额外问题。
另一方面，要为非紧定义域建立弗雷德霍姆理论，则需要复杂得多的分析。（此时常常研究
扰动方程；比较 \srcite{116} 中弗洛尔理论的发展。）这两种推广都超出本书范围，而要理解
\BookRecordIndex{T0143}格罗莫夫--威滕不变量的完整结构，两者又都不可缺少；例如可参见 Eliashberg、Givental 和
Hofer 最近关于\BookRecordIndex{T0352}辛场论的论文 \srcite{101} 中所展示的结构。

\paragraph{紧致性。}
下一项任务是弄清楚\BookRecordIndex{T0246}模空间 $\mathcal M^*(A;J)$ 何时紧致。在这里，\BookRecordIndex{T0372}驯服条件起着本质作用。
\BookRecordIndex{T0363}辛形式 $\omega$ 与一个被 $\omega$ \BookRecordIndex{T0372}驯服的\BookRecordIndex{T0088}殆复结构 $J$ 一起确定 $M$ 上的黎曼度量
\[
  \langle v,w\rangle=\frac12\bigl(\omega(v,Jw)+\omega(w,Jv)\bigr),
\]
而相对于这一度量，\BookRecordIndex{T0003}$J$-全纯曲线 $u:\Sigma\to M$ 的能量由下式给出：
\[
  E(u)=\frac12\int_\Sigma |du(z)|_J^2\,d\operatorname{vol}_\Sigma
      =\int_\Sigma u^*\omega.
\]
因此，\BookRecordIndex{T0003}$J$-全纯曲线微分的 $L^2$-范数有一个只依赖于 $u$ 所表示同调类 $A$ 的一致上界。
这本身不能保证紧致性，因为它恰处在索伯列夫估计的临界情形。（若 $p>2$ 时 $du$ 的
$L^p$-范数有一致上界，就能保证紧致性。）

紧致性失效的另一种表现，是能量 $E(u)$ 在 $\Sigma$ 上度量的共形缩放下不变。当曲线的
定义域 $\Sigma$ 是黎曼球面 $\mathbb{CP}^1$ 时，这一现象尤其明显，因为此时存在一大群
整体的、而非仅仅局部的缩放。事实上，非紧群 $G=\operatorname{PSL}(2,\mathbb C)$ 通过
共形变换
\[
  z\longmapsto\frac{az+b}{cz+d}
\]
作用在黎曼球面上。因此每个 $u\in\mathcal M^*(A;J)$ 都有一个非紧的重参数化族
$u\circ\phi$，其中 $\phi\in G$；所以 $\mathcal M^*(A;J)$ 本身不可能既紧致又非空
（除非 $A$ 是零类，此时所有映射 $u$ 都是常值映射）。不过，商空间
$\mathcal M^*(A;J)/G$ 有时是紧致的。

回忆：若同调类 $B\in H_2(M)$ 位于胡雷维茨（Hurewicz）同态
$\pi_2(M)\to H_2(M)$ 的像中，就称 $B$ 为\BookRecordIndex{T0285}球面类。下面是第 4 章第一个重要定理的陈述。
注意，对 $A$ 的假设蕴含任何 $A$-曲线都不可能是多重覆盖的，即
$\mathcal M^*(A;J)=\mathcal M(A;J)$。

\begin{namedtheorem}[B]
假设不存在球面同调类 $B\in H_2(M)$ 使得
$0<\omega(B)<\omega(A)$。那么\BookRecordIndex{T0246}模空间 $\mathcal M^*(A;J)/G$ 是紧致的。
\end{namedtheorem}

为证明这一点，先说明：如果序列 $u_\nu\in\mathcal M^*(A;J)$ 在
$\mathcal M^*(A;J)$ 中没有聚点，就存在 $z\in\mathbb{CP}^1$，使微分
$du_\nu(z)$ 无界。取子列后，可找到 $z$ 在 $\mathbb{CP}^1$ 中的一列递减邻域 $U_\nu$，
使其像 $u_\nu(U_\nu)$ 在极限中收敛到一个非常值 $J$-全纯球面。若该球面表示同调类 $B$，
则或者 $\omega(B)=\omega(A)$，此时可以对 $u_\nu$ 重参数化，使其在
$\mathcal M^*(A;J)$ 中收敛；或者 $\omega(B)$ 严格介于 $0$ 与 $\omega(A)$ 之间。
这就是“\BookRecordIndex{T0243}冒泡”过程；在当前情形中，它以一种简单而且几何上清楚的方式发生，参见第 4.2 节。

定理 B 蕴含：若 $\omega(A)$ 已经是 $\omega$ 在球面上所取的最小正值，则\BookRecordIndex{T0246}模空间
$\mathcal M^*(A;J)/G$ 紧致。如下文所见，这已经足以导出若干有趣应用。不过，为处理更
一般的类 $A$，我们必须理解序列 $u_\nu$ 的完整极限。解决这一困难问题所需的分析在
第 4 章后半部发展。特别地，为了看出极限曲线连通，我们在第 4.7 节研究小能量长圆柱的性质。

极限的完整结构由定理 5.3.1 表述，这一结果称为格罗莫夫紧致性定理。我们采用
Kontsevich 的\BookRecordIndex{T0340}稳定映射语言，给出比格罗莫夫在 \srcite{160} 中原始陈述更现代的形式。
尽管本书的应用并不严格需要这种语言，它却使我们能够在自然框架中建立\BookRecordIndex{T0143}格罗莫夫--威滕
不变量，而且是该理论近期许多应用、特别是与\BookRecordIndex{T0229}量子上同调有关的应用之基础。因此我们在
第 5 章和附录 D 中相当详细地发展这一语言。

初学者最合适的策略，或许是先设法理解\BookRecordIndex{T0141}格罗莫夫收敛的定义（定义 5.2.1）和定理 5.3.1
的陈述，把更多细节留到以后。命题 4.1.5 也很重要：它断言，对每个 $J$ 和每个 $c>0$，
具有 $J$-全纯代表且能量 $\omega(A)$ 不超过 $c$ 的类 $A$ 只有有限多个。
